% =============================================================================
% Kapitel 2: Theoretische Grundlagen
% =============================================================================
\chapter{Theoretische Grundlagen}
\label{ch:theorie}

\section{Röntgenabsorption und die Absorptionskante}

Trifft ein monochromatischer Röntgenstrahl der Intensität $I_0$ auf eine
Probe der Dicke $d$, so wird die transmittierte Intensität durch das
Lambert-Beersche Gesetz beschrieben:
\begin{equation}
  I(E) = I_0\,\mathrm{e}^{-\mu(E)\,d}\,,
  \label{eq:lambert}
\end{equation}
wobei $\mu(E)$ der lineare Absorptionskoeffizient ist.
Wird die Photonenenergie $E$ über eine Ionisierungsschwelle eines in der
Probe enthaltenen Elements hinweg erhöht, steigt $\mu(E)$ sprunghaft
an -- man spricht von einer \emph{Absorptionskante}.
An der K-Kante beispielsweise wird ein 1s-Elektron in das Kontinuum
angeregt; die Kantenenergie $E_0$ ist elementspezifisch und liegt für
die 3d-Übergangsmetalle im Bereich von ca.\ \SIrange{5}{9}{\kilo\electronvolt}
\cite{Koningsberger1988}.

Die Absorptionskante ist nicht nur ein Schwellenwert, sondern enthält in
ihrer näheren Umgebung eine reiche Struktur:
Der Bereich bis etwa \SI{30}{\electronvolt} oberhalb von $E_0$ wird
als XANES (X-ray Absorption Near Edge Structure) bezeichnet und
spiegelt die elektronische Struktur und Symmetrie des Absorbers wider.
Der sich anschließende Bereich bis zu einigen hundert Elektronenvolt
oberhalb der Kante heißt EXAFS (Extended X-ray Absorption Fine
Structure) und wird durch die Interferenz des emittierten
Photoelektrons mit seinen an den Nachbaratomen gestreuten Anteilen
verursacht.

% XAS-Prinzipskizze (TikZ)
% Standalone TikZ figure: XAS measurement principle
% Can be compiled separately or included via \input
\begin{figure}[htbp]
  \centering
  \begin{tikzpicture}[
    scale=0.9,
    atom/.style={circle, draw, fill=blue!20, minimum size=0.5cm,
                 inner sep=0pt, font=\tiny},
    absorber/.style={circle, draw, fill=red!40, minimum size=0.6cm,
                     inner sep=0pt, font=\tiny\bfseries},
    photon/.style={-{Stealth}, decorate,
                   decoration={snake, amplitude=1.5pt, segment length=6pt},
                   thick, red!70!black},
    electron/.style={-{Stealth}, thick, blue!70!black, dashed},
    scattered/.style={-{Stealth}, thick, green!50!black, dashed},
  ]
    % X-ray beam
    \draw[photon] (-3, 0) -- (-0.5, 0)
      node[midway, above=4pt, font=\small, text=red!70!black] {$h\nu$};

    % Absorber atom
    \node[absorber] (abs) at (0,0) {A};

    % Neighbor atoms
    \node[atom] (n1) at (2.2, 0.8) {N};
    \node[atom] (n2) at (2.0, -1.0) {N};
    \node[atom] (n3) at (-1.0, 2.0) {N};
    \node[atom] (n4) at (1.5, 2.2) {N};

    % Photoelectron wave (outgoing)
    \draw[electron] (abs.east) -- (n1.west)
      node[midway, above, font=\scriptsize] {$e^-$};
    \draw[electron] (abs.east) -- (n2.north west);

    % Scattered wave (returning)
    \draw[scattered] (n1.south west) -- (abs.north east);
    \draw[scattered] (n2.north) -- (abs.south east);

    % Interference annotation
    \draw[<->, thick, gray] (0, -2) -- (2.2, -2)
      node[midway, below, font=\small] {$R_j$};

    % Labels
    \node[font=\small, text=red!70!black, below left] at (-0.3, -0.5)
      {Absorber};
    \node[font=\small, text=blue!70!black, right] at (2.5, 0.8)
      {Nachbar};

    % Right side: schematic spectrum
    \begin{scope}[xshift=7cm, yshift=-0.5cm]
      \draw[->, thick] (0, 0) -- (4.5, 0)
        node[right, font=\small] {$E$};
      \draw[->, thick] (0, 0) -- (0, 3.5)
        node[above, font=\small] {$\mu(E)$};

      % Pre-edge
      \draw[thick, blue!60!black] (0.2, 0.5) -- (1.2, 0.6);

      % Edge jump
      \draw[thick, blue!60!black] (1.2, 0.6) -- (1.4, 2.5);

      % EXAFS oscillations
      \draw[thick, blue!60!black]
        (1.4, 2.5) .. controls (1.7, 2.8) and (1.9, 2.2) ..
        (2.1, 2.5) .. controls (2.3, 2.8) and (2.5, 2.3) ..
        (2.7, 2.5) .. controls (2.9, 2.7) and (3.1, 2.4) ..
        (3.3, 2.5) .. controls (3.5, 2.6) and (3.7, 2.45) ..
        (4.0, 2.5);

      % Annotations
      \draw[<->, gray, thick] (1.15, 0.6) -- (1.15, 2.5)
        node[midway, left, font=\scriptsize] {$\Delta\mu_0$};
      \node[font=\scriptsize] at (1.3, 0.2) {$E_0$};
      \draw[gray, dashed] (1.3, 0) -- (1.3, 0.6);

      \draw[decorate, decoration={brace, amplitude=4pt}]
        (1.5, 3.2) -- (4.0, 3.2)
        node[midway, above=4pt, font=\scriptsize] {EXAFS};

      \draw[decorate, decoration={brace, amplitude=4pt}]
        (1.1, 3.2) -- (1.5, 3.2)
        node[midway, above=4pt, font=\scriptsize] {XANES};
    \end{scope}
  \end{tikzpicture}
  \caption{Prinzip der Röntgenabsorptionsspektroskopie.
    Links: Das einfallende Photon erzeugt am Absorberatom (A) ein
    Photoelektron, das an den Nachbaratomen (N) gestreut wird.
    Die Interferenz zwischen ausgehender und rückgestreuter Welle
    moduliert den Absorptionskoeffizienten.
    Rechts: Schematisches Absorptionsspektrum $\mu(E)$ mit
    Kantensprung $\Delta\mu_0$, XANES- und EXAFS-Bereich.}
  \label{fig:xas-principle}
\end{figure}


\section{Die EXAFS-Gleichung}
\label{sec:exafs-gleichung}

Die EXAFS-Oszillation wird als normierter, kantenbereinigte Anteil des
Absorptionskoeffizienten definiert:
\begin{equation}
  \chi(E) = \frac{\mu(E) - \mu_0(E)}{\Delta\mu_0(E_0)}\,,
  \label{eq:chi_E}
\end{equation}
mit dem atomaren Hintergrund $\mu_0(E)$ und dem Kantensprung
$\Delta\mu_0(E_0)$.
Die Umrechnung in den Wellenzahlraum erfolgt über
\begin{equation}
  k = \sqrt{\frac{2m_e(E - E_0)}{\hbar^2}}\,,
  \label{eq:k_def}
\end{equation}
sodass das EXAFS-Signal als Funktion von $k$ geschrieben werden kann.

Im Einfachstreumodell ergibt sich die bekannte EXAFS-Gleichung
\cite{Stern1974, Sayers1971}:
\begin{equation}
  \chi(k) = S_0^2 \sum_j \frac{N_j\,|f_j(k)|}{k\,R_j^2}\,
    \sin\!\bigl(2kR_j + \delta_j(k)\bigr)\,
    \mathrm{e}^{-2\sigma_j^2 k^2}\,
    \mathrm{e}^{-2R_j/\lambda(k)}\,.
  \label{eq:exafs}
\end{equation}

Die Summation läuft über alle Koordinationsschalen $j$ mit folgenden Größen:
\begin{itemize}[nosep]
  \item $N_j$ -- Koordinationszahl (Anzahl der Nachbarn in Schale $j$),
  \item $R_j$ -- mittlerer Abstand zum Absorber,
  \item $f_j(k)$ und $\delta_j(k)$ -- Rückstreuamplitude und
        Phasenverschiebung, die vom Streuerelement abhängen,
  \item $\sigma_j^2$ -- mittlere quadratische Abweichung des Abstands
        (Debye-Waller-Faktor), die sowohl thermische Schwingungen als auch
        statische Unordnung erfasst,
  \item $\lambda(k)$ -- mittlere freie Weglänge des Photoelektrons,
  \item $S_0^2$ -- Amplitudenreduktionsfaktor, der inelastische
        Vielteilcheneffekte berücksichtigt.
\end{itemize}

Diese Gleichung verdeutlicht den zentralen Informationsgehalt eines
EXAFS-Spektrums: Jede Nachbarschale erzeugt eine Sinusschwingung im
$k$-Raum, deren Frequenz vom Abstand $R_j$ und deren Amplitude von der
Koordinationszahl $N_j$, der Rückstreuamplitude und dem
Debye-Waller-Faktor abhängt.

\section{Fouriertransformation und $R$-Raum}

Durch Fouriertransformation von $\chi(k)$ in den $R$-Raum lassen sich
die Beiträge verschiedener Koordinationsschalen separieren.
Üblicherweise wird das $k$-gewichtete Signal $\chi(k)\cdot k^n$ mit
einer Fensterfunktion $W(k)$ multipliziert:
\begin{equation}
  \tilde{\chi}(R) = \frac{1}{\sqrt{2\pi}}
    \int_{k_{\min}}^{k_{\max}} \chi(k)\,k^n\,W(k)\,
    \mathrm{e}^{2ikR}\,\mathrm{d}k\,.
  \label{eq:ft}
\end{equation}
Der Betrag $|\tilde{\chi}(R)|$ zeigt Maxima bei Abständen, die in erster
Näherung den Positionen der Koordinationsschalen entsprechen.
Zu beachten ist, dass die Peaks gegenüber den tatsächlichen
Bindungslängen um etwa
$\SIrange{0.2}{0.5}{\Angstrom}$ zu kleineren Werten verschoben sind,
da die Phasenverschiebung $\delta_j(k)$ nicht berücksichtigt wird.

Die Wahl der $k$-Gewichtung beeinflusst die relative Empfindlichkeit:
$k^1$ betont leichte Rückstreuer bei kleinen $k$-Werten,
$k^3$ hebt schwere Streuer bei hohen $k$ hervor,
und $k^2$ stellt einen häufig gewählten Kompromiss dar \cite{Bunker2010}.

\section{Mehrfachstreuung}

Gleichung~\eqref{eq:exafs} beschreibt ausschließlich die Einfachstreuung
(single scattering, SS), bei der das Photoelektron auf direktem Weg zum
Nachbaratom und zurück streut.
Bei kollinearen oder nahezu kollinearen Atomanordnungen können jedoch
Mehrfachstreupfade (multiple scattering, MS) einen erheblichen Beitrag
leisten.
Insbesondere Drei- und Vierbeinpfade, bei denen das Elektron an
Zwischenatomen umgelenkt wird, sind in dicht gepackten Strukturen wie
FCC-Metallen von Bedeutung, da dort die Fokussierung entlang der
$\langle 110\rangle$-Richtung den Effekt verstärkt \cite{Zabinsky1995}.

Die Software \feff{} \cite{Rehr2010} berechnet sowohl Einfach- als auch
Mehrfachstreupfade ab initio unter Verwendung der
Muffin-Tin-Approximation für das Kristallpotential und einer
selbstkonsistenten Berechnung der Streuphasen.
Die Berücksichtigung von MS-Pfaden ist insbesondere bei der Analyse
höherer Koordinationsschalen unerlässlich.

\section{Wavelet-Transformation}
\label{sec:wavelet}

Eine konventionelle Fouriertransformation liefert nur $R$-aufgelöste
Information; die Zuordnung eines Peaks zu einem bestimmten Streuerelement
ist allein aus $|\tilde{\chi}(R)|$ nicht möglich, da die $k$-Abhängigkeit
verloren geht.
Die Wavelet-Transformation (WT) umgeht diese Einschränkung, indem sie
eine zweidimensionale Darstellung in $(k,R)$ erzeugt
\cite{Funke2005, Timoshenko2009}.
Dadurch lassen sich Streubeiträge bei gleichem Abstand, aber
unterschiedlichen Rückstreuelementen (und damit unterschiedlicher
$k$-Abhängigkeit von $f_j(k)$) voneinander trennen.

In \evax{} wird der Vergleich zwischen Theorie und Experiment wahlweise
im $k$-Raum, im $R$-Raum oder im Wavelet-Raum durchgeführt.
Der Wavelet-Raum stellt dabei die schärfste Bedingung an die
Übereinstimmung, da er beide Dimensionen gleichzeitig berücksichtigt
\cite{Timoshenko2014}.

\section{Chemische Nahordnung: der Warren-Cowley-Parameter}
\label{sec:sro}

In einer Legierung mit mehreren Komponenten interessiert nicht nur die
mittlere Struktur, sondern auch die Frage, ob bestimmte
Elementkombinationen als Nachbarn bevorzugt oder vermieden werden.
Der Warren-Cowley-Parameter quantifiziert diese chemische Nahordnung
(short-range order, SRO) \cite{Warren1969, Cowley1950}:
\begin{equation}
  \alpha_{ij} = 1 - \frac{P_{ij}}{c_j}\,,
  \label{eq:warren_cowley}
\end{equation}
wobei $P_{ij}$ die bedingte Wahrscheinlichkeit ist, in der nächsten
Nachbarschale eines $i$-Atoms ein $j$-Atom zu finden, und $c_j$ der
Molenbruch von Element $j$ ist.

Die Interpretation ist unmittelbar anschaulich:
\begin{itemize}[nosep]
  \item $\alpha_{ij} = 0$: statistisch zufällige Verteilung (ideale Mischkristall),
  \item $\alpha_{ij} > 0$: $j$-Nachbarn werden in der Umgebung von $i$ vermieden,
  \item $\alpha_{ij} < 0$: $j$-Nachbarn werden bevorzugt.
\end{itemize}

Da \evax{} ein vollständiges atomistisches Modell liefert, lassen sich
$P_{ij}$ und damit $\alpha_{ij}$ direkt aus der finalen Struktur
berechnen, ohne zusätzliche Annahmen über die Form der
Paarverteilungsfunktion.

\section{Grenzen der EXAFS-Methode}

Trotz ihrer Stärken unterliegt die EXAFS-Analyse einigen
grundsätzlichen Einschränkungen:
\begin{itemize}
  \item \textbf{Lokale Sonde:} EXAFS erfasst nur die gemittelte
        Nahordnung innerhalb weniger Ångström um das Absorberatom.
        Langreichweitige Ordnung (Domänengrößen, Gitterkonstanten im
        klassischen Sinn) ist der Methode nicht zugänglich.
  \item \textbf{Parameterkorrelation:} $N_j$ und $\sigma_j^2$ sind
        stark korreliert, was die gleichzeitige Bestimmung beider Größen
        aus einem einzelnen Spektrum erschwert \cite{Bunker2010}.
  \item \textbf{Begrenzte Informationsdichte:} Die Nyquist-Zahl
        (Gl.~\eqref{eq:nyquist}) limitiert die Zahl der zuverlässig
        bestimmbaren Parameter.
  \item \textbf{Mehrkomponentige Systeme:} Bei mehr als zwei
        Streuerelementen in der gleichen Schale werden die Beiträge
        zunehmend schwer trennbar, sofern nicht mehrere Kanten
        gleichzeitig ausgewertet werden.
\end{itemize}

Gerade der letzte Punkt motiviert den Einsatz von Methoden wie \evax{},
die durch simultane Multi-Edge-Auswertung die Informationsdichte
erheblich steigern.
